% Lösungen Prozentrechnung · Umwandeln · Prüfungs-Fokus MSA (zusammenbau v0.6)
\einheitenkopf{Einheit 1 · Prozente als Anteile}

\erg{1}{a) $50\,\%$ \quad b) $37\,\%$}
\erg{2}{a) $5$ von $100$ Paketen kommen zu spät. ($5\,\% = \frac{5}{100}$; „jedes 5.“ wären $20\,\%$) \quad b) $8$ von $100$ Gästen bestellen Tee. ($8\,\% = \frac{8}{100}$; „jeder 8.“ wären $12{,}5\,\%$) \quad c) $25\,\%$ ($\frac{1}{4} = \frac{25}{100}$) \quad d) $5\,\%$ ($\frac{1}{20} = \frac{5}{100}$)}
\erg{3}{a) Aussage 2 ist falsch: $200 : 600 \approx 0{,}33$, also ein Drittel. Richtig: Ein Drittel der Kinder fährt mit dem Rad. (Aussage 1 stimmt: $150 : 600 = 0{,}25$.) \quad b) Aussage 1 ist falsch: $600 : 2\,400 = 0{,}25$, also ein Viertel. Richtig: Jedes vierte Mitglied ist ein Kind. (Jedes fünfte, $480 : 2\,400 = 0{,}2$, gilt für die Jugendlichen; Aussage 2 stimmt: $1\,200 : 2\,400 = 0{,}5$.) \quad c) Aussage 1 ist falsch: auf dem Land sind es $\frac{1}{10} = 10\,\%$. Richtig: Ein Zehntel der Befragten auf dem Land geht in die Bibliothek (ein Viertel, $\frac{1}{4} = 25\,\%$, gilt für die Stadt). Aussage 2 stimmt: $55\,\% > 50\,\%$. \quad d) Aussage 2 ist falsch: $\frac{1}{20} = 5\,\%$. Richtig: Jeder zwanzigste Grundschüler frühstückt nie zu Hause (jeder zehnte gilt für die Oberschule). Aussage 1 stimmt: $\frac{3}{5} = 60\,\%$. \quad e) Nein: $27{,}5\,\% + 22{,}0\,\% = 49{,}5\,\%$, das ist weniger als die Hälfte. \quad f) Ja: $45{,}3\,\% + 21{,}8\,\% + 4{,}3\,\% = 71{,}4\,\%$ (oder $100\,\% - 28{,}6\,\%$), das ist mehr als $70\,\%$.}
\erg{4}{a) Biogas: $100 - 38 - 24 - 9 - 3 = 26$, also $26\,\%$ \quad b) Siedlung: $100 - 34 - 41 - 4 - 2 = 19$, also $19\,\%$}
